\section{Introduction}\label{sec:introduction}

Let $X$ be a normal projective klt variety with $K_X\sim_\Q0$.
Its canonical index is the least positive integer $r$ for which
$rK_X$ is an integral principal divisor. The index problem asks for
a bound depending only on the dimension. More generally, for a log
Calabi--Yau pair one allows dependence on the boundary coefficients.
In dimension three, Jiang obtains a common canonical multiple for
all klt Calabi--Yau threefolds, using boundedness modulo flops to
handle the noncanonical case \cite[Theorem~1.6 and Corollary~1.7]{Jiang2021}.
Xu's inductive results give the lc threefold index theorem and the
klt fourfold theorem with nonzero boundary
\cite[Theorems~1.13--1.14]{Xu2019}. Jiang and Haidong Liu bound
surface log pluricanonical representations and, through Fujino's
gluing framework, obtain the index theorem for slc pairs through
dimension three and connected non-klt lc fourfold pairs
\cite[Theorem~1.4 and Corollaries~1.6--1.7]{JiangLiu2021}.
Birkar also proves an index bound in every fixed dimension for
rationally connected klt log Calabi--Yau pairs with coefficients in
a fixed rational set satisfying the descending chain condition
(DCC) \cite[Corollary~1.8]{Birkar2025}.
The distinction between zero and nonzero boundary, and between
rationally connected and non-uniruled varieties, is therefore central
to the remaining fourfold argument.

\begin{theorem}[Uniform canonical index]\label{thm:fourfold-index-zero}
There is a positive integer $N_4$ such that $N_4K_X\sim0$ for every
normal projective klt complex fourfold $X$ with $K_X\sim_\Q0$.
\end{theorem}

The conclusion concerns integral Weil divisors: it gives a
trivialization in one degree for all the fourfolds in the statement.
Combined with the preceding non-klt and nonzero-boundary results,
it gives the lc fourfold index theorem for any fixed finite rational
coefficient set. In applying Xu's theorem, a finite set is included
in the hyperstandard set generated by its complementary coefficients.
No assertion for arbitrary real coefficient sets is intended.

Our effective application concerns the map determined by a log
canonical divisor. For a rational Cartier divisor $D$ on a normal
integral projective variety, fix a divisor representative and put
\[
 V_m(D)=H^0\bigl(X,\OO_X(\lfloor mD\rfloor)\bigr),\qquad m\geq1.
\]
The sheaf here is the rank-one reflexive divisorial sheaf, so this
definition does not require $mD$ to be Cartier. The Iitaka field
$K(D)\subset\C(X)$ is generated by ratios of nonzero sections in
the same $V_m(D)$, over all $m$. A complete system gives the Iitaka
fibration when its ratios generate this entire field. Merely obtaining
the correct image dimension would not rule out a finite extension.

\begin{theorem}[Uniform effective log Iitaka fibrations]
\label{thm:main}
Let $\Phi\subset[0,1]\cap\Q$ be finite. There is a positive integer
$m=m(\Phi)$ such that the following holds. Suppose $X$ is a normal
integral projective complex fourfold, $\Delta\geq0$ is a rational
Weil divisor whose nonzero coefficients lie in $\Phi$, $(X,\Delta)$
is lc, and $D=K_X+\Delta$ is rational Cartier and pseudo-effective.
Then the complete system $|\lfloor mD\rfloor|$ is nonempty and its
section ratios generate $K(D)$ inside $\C(X)$. The same integer works
for every choice of canonical divisor on $X$.
\end{theorem}

In Kodaira dimension four the system is birational; in Kodaira
dimension zero it is nonempty with image a point. Nonvanishing follows
from pseudo-effectivity by the good-model theorem and effective
exceptional comparison in OpenAI, \emph{Log abundance in characteristic
zero} \cite[Theorems~1.1 and~11.1, Lemma~2.2]{OpenAIAbundance}.
This qualitative step supplies a section in some positive degree,
not a uniform degree; the uniform $m(\Phi)$ comes from the index and
fibration arguments below.

Earlier effective results separate the big case from lower-dimensional
Iitaka bases. Hacon--McKernan--Xu give a uniform birational degree
for big lc adjoints with coefficients in any fixed DCC set
\cite[Theorem~1.3]{HMX2014}. Viehweg--Zhang treat smooth varieties
of Kodaira dimension two. Their degree depends on invariants of
the general fibre; for fourfolds the fibre is a surface of Kodaira
dimension zero, and boundedness of those invariants makes the degree
universal \cite[Theorem~0.2 and the discussion preceding Corollary~0.4]{ViehwegZhang2009}.
Todorov--Xu treat klt pairs of log Kodaira codimension two in
dimensions two, three and four, with coefficients in a rational DCC set
\cite[Theorem~1.2]{TodorovXu2009}. For fourfolds their hypothesis is
$\kappa(K_X+\Delta)=2$. Chen, Han and Jihao Liu formulate the
effective log Iitaka conjecture for lc pairs of nonnegative Kodaira
dimension with coefficients in an arbitrary DCC subset of $[0,1]$,
and prove it through dimension three
\cite[Conjecture~1.2 and Corollary~1.5]{ChenHanLiu2023}.

The canonical bundle formula relates this question to effective
birationality on the base. Birkar--Zhang prove effective Iitaka
fibrations for smooth projective varieties of nonnegative Kodaira
dimension, with constants depending on the dimension, the least
nonvanishing pluricanonical degree on the general fibre, and the
middle Betti number of a smooth model of its canonical cover.
In fixed dimension and
for a fixed DCC coefficient set, their polarized effective
birationality theorem gives a uniform degree for a big base adjoint
once the nef part has bounded denominator
\cite[Theorems~1.2--1.3]{BirkarZhang2016}.
Theorem~\ref{thm:main} gives the fourfold conclusion for finite
rational coefficient sets, across all nonnegative Kodaira dimensions
and for lc singularities. Its coefficient scope is narrower than
that of the DCC conjecture and the earlier big and klt results.
The uniform degree uses the canonical-index theorem and the
relative denominator theorem below. The passage from
pseudo-effectivity to nonvanishing is the separate qualitative
input described above.

\subsection{The model comparison and the index reductions}

The uniform argument needs a model on which the adjoint defines a morphism. The
comparison in Proposition~\ref{prop:semiample-model} preserves
\emph{every} space $V_m(D)$ under a crepant dlt modification and an
adjoint-negative program. Its semiample output uses the
abundance-after-nonvanishing theorem of OpenAI,
\emph{Lifting sections from the reduced support of an adjoint}
\cite[Theorem~\BANumber]{OpenAILifting}. That theorem
applies to projective lc rational pairs over $\C$ of dimension at
most four: a nef rational Cartier adjoint of nonnegative Kodaira
dimension is semiample. In Kodaira dimension zero it is rationally
linearly trivial. We state the model comparison before the index
reduction because it is needed there as well as in the final
effective application.

The singular Beauville--Bogomolov decomposition gives a finite
quasi-\'etale cover by a product of an abelian variety and nonflat
Calabi--Yau or symplectic factors. Druel proves the needed
dimension-at-most-five result \cite[Theorem~1.2]{Druel2018};
H\"oring--Peternell prove the arbitrary-dimensional version
\cite[Theorem~1.5]{HoringPeternell2019}, using Druel's integrability
results and the holonomy decomposition of Greb--Guenancia--Kebekus
\cite[Theorem~B and Proposition~D]{GGK2019}.
Here quasi-\'etale means \'etale outside a subset of codimension
at least two. Integral cohomology bounds the abelian characters,
and lower-dimensional pluricanonical characters handle products.
The remaining single-factor case has rationally connected proper
rational images. Its noncanonical quotients are handled by a
relative torsion bound; canonical quotients admit a crepant
terminal model $V$ with the same canonical index, irregularity
zero and countable birational self-map group.

An effective divisor $D$ of Iitaka dimension one, two or three on
$V$ would give another bounded-index case. A small rational
boundary $\delta D$ produces a semiample contraction, but the
relative index theorem is applied only after proving that the
underlying canonical divisors remain crepant. The resulting
fibration has boundary zero, so the relative denominator and
rationally connected torsion theorems
\cite[Theorem~\RDenominatorNumber{} and Proposition~\RCTorsionNumber]{OpenAIRelative}
have constants independent of $D$ and $\delta$.
Consequently, in the unbounded-index case every effective divisor
of positive Iitaka dimension is big. This implies the precise
fibration property used in the scalar estimate: every intermediate
equivariant rational fibration of the canonical cover has a
geometric generic fibre of general type on a smooth projective
resolution. It is a conclusion of the reduction.

\subsection{Curves on the canonical cover and its products}

Let $r$ be the canonical index of the remaining terminal fourfold
$V$. A rational function trivializing $rK_V$ defines its canonical
cyclic cover $\pi:Y\to V$, a connected degree-$r$ quasi-\'etale cover
with deck group $\mu_r$ and trivial canonical class upstairs. The
minimality of $r$ is what makes this the primitive cover. Choose an
ample rational divisor $L$ on $V$, scaled so that its volume is
bounded above and below by absolute positive constants. Then
$(\pi^*L)^4=rL^4$.

For an ample rational divisor $P$, two local measurements will be
used. The number $\gamma(P;V)$ is the supremum of the order of
vanishing of a section in degree $k$, divided by $k$, at a very
general smooth point; degrees are sufficiently divisible to make
the divisor Cartier. Its precise definition is given in
Definition~\ref{def:orders}. For an integral curve $C$ through a
smooth marked point $x$, the ratio
\[
 \frac{P\cdot C}{\operatorname{mult}_x C}
\]
compares its degree with its multiplicity at that point. The infimum
over such curves is the Seshadri constant $\varepsilon(P;x)$.
The scalar estimates bound $\gamma$ above and $\varepsilon$ below
by absolute multiples of the fourth root of the volume, under the
fibration hypothesis just obtained. Applied upstairs, the lower
bound makes every such curve ratio grow at least as $r^{1/4}$.

Now take $t$ copies of $Y$ and divide by the simultaneous deck action:
\[
 Z_t=Y^t/\mu_{r,\mathrm{diag}},\qquad
 \theta_t:Z_t\longrightarrow V^t.
\]
The polarization $P_t$ is the pullback by $\theta_t$ of the sum of
the $t$ copies of $L$. On $Z_t$, the product estimate gives curves
through very general points whose degree divided by marked
multiplicity is bounded by a constant times $t^2$, independently of
$r$. This estimate concerns the quotient by the diagonal group;
it does not assert the same bound on $Y^t$.
The maps are displayed in Figure~\ref{fig:diagonal-cover}.

\begin{figure}[ht]
\centering
\begin{tikzcd}[column sep=large,row sep=large]
Y^t \arrow[r,"/\mu_{r,\mathrm{diag}}"] \arrow[d,"\operatorname{pr}_1"']
 & Z_t \arrow[r,"\theta_t"] \arrow[d,"q_1"]
 & V^t \arrow[d,"\operatorname{pr}_1"]\\
Y \arrow[r,"\pi"'] & V \arrow[r,equal] & V
\end{tikzcd}
\caption{The product polarization is pulled back through $\theta_t$.
The geometric generic fibre of $q_1$ is $Y^{t-1}$: choosing one lift of
the marked point in $V$ removes the simultaneous deck action. This
identification allows the growing scalar curve bound upstairs to be
tested inside a product whose curve cost is bounded independently of $r$.}
\label{fig:diagonal-cover}
\end{figure}


The product estimate is proved by reducing a fixed set of data to
positive characteristic and comparing ordinary jets with Frobenius
jets. The local comparison is that of Musta\c{t}\u{a}--Schwede
\cite[Equation (2.8)]{MustataSchwede2014}. The canonical filtration described by Sun
\cite[Theorem 3.7 and Lemma 4.2]{Sun2008}, in the ideal-theoretic
description of Kitadai--Sumihiro
\cite[Definition 3.1 and Remark 3.2]{KitadaiSumihiro2008}, supplies
the truncated symmetric powers and their determinant factors.
We prove the slope estimates and the diagonal rank comparison
needed for these products. For each variety, cover and product size,
the arithmetic model is fixed before the characteristic tends to
infinity. The later contradiction fixes a finite list of product
sizes before allowing $r$ to grow. These are separate limits;
neither requires a uniform arithmetic model for all $r$.

\subsection{From chains of curves to birational maps}

The curve estimates concern different spaces. We compare them by
forming chains from finitely many covering families of marked
curves on $Z_t$, chosen compatibly with coordinate projections and
the group actions. Each movement follows a chosen curve from its
marked point to a generic point. The closures of the generically reachable sets
are fibres of a rational quotient, called \emph{chain leaves}.
The construction uses Campana's stabilization strategy
\cite[Theorem~1.1 and its proof sketch]{Campana2004};
parameter differentiation in the manner of Ein--K\"uchle--Lazarsfeld
\cite[Proposition~2.3]{EinKuchleLazarsfeld1995} retains a degree
bound. In our setting a generic point of an $e$-dimensional leaf
can be reached in at most $e$ movements. If the marked curve
ratios are at most $B$, the leaf degree is at most $(eB)^e$.
We prove the projection and finite-quotient compatibilities needed
to use these bounds together.

Write $H_t\subset Z_t$ for a generic leaf and
$W_t=\theta_t(H_t)\subset V^t$. A positive-dimensional fibre of
$H_t\to V$ would lie, after choosing one lift of the marked base
point, in $Y^{t-1}$. The scalar lower bound there grows with $r$,
contradicting the bounded leaf degree. Thus, for sufficiently large
$r$, the maps to each coordinate are generically finite onto their
images and the leaf dimensions lie between one and four. Among a
sufficiently long finite list of product sizes, stabilization of
dimensions and a comparison of forgetting-map degrees force one
map $W_{t+1}\to W_t$ to have degree one.

This birational forgetting map produces more than a dimension
contradiction. The leaf tangent spaces identify positive-dimensional
subspaces $A_x\subset T_xV$ at different points. We show that the
resulting linear transports depend only on their two endpoints
and compose. A vector transported from one fixed point then defines
a rational vector field. Degree-one forgetting supplies rational
evaluation of its local flow, and inverse flow gives birational
self-maps. The resulting uncountable family contradicts the
countability of $\Bir(V)$, completing the index argument.

Finally, the model comparison turns the effective theorem into a
question on the base of a semiample contraction. For every
positive-dimensional base the relative big-base theorem
\cite[Proposition~\REffectiveNumber]{OpenAIRelative} gives the
whole base field in a uniform complete floor system, including when
the contraction is birational. In
Kodaira dimension zero, the three index cases give a nonzero section
in one uniform degree. The comparison of all
integer degrees brings these systems back to the original pair.
The relative results used in this argument are proved independently
of Theorem~\ref{thm:fourfold-index-zero}. The abundance-after-nonvanishing
input from \cite{OpenAILifting} and the uniform and index arguments above
do not require logarithmic or orbifold subadditivity. The full
pseudo-effective statement of Theorem~\ref{thm:main}, however, first uses
\cite{OpenAIAbundance} to obtain nonvanishing and inherits that companion's
logarithmic Iitaka subadditivity input.

\subsection{Organization and conventions}

Section~\ref{sec:reductions} constructs the semiample model and compares
all rounded systems. Section~\ref{sec:relative-inputs} records the exact
relative inputs. Section~\ref{sec:index-reductions} removes product and
rationally connected cases and proves the fibration condition needed by
the scalar estimate. Sections~\ref{sec:chains}--\ref{sec:transports}
prove the chain, scalar, product and transport statements in that order.
Section~\ref{sec:index-zero} applies them to the remaining canonical
covers. Section~\ref{sec:main-proof} derives the effective theorem.

Varieties are over $\C$ unless a geometric generic fibre or an
arithmetic reduction is specified. Canonical divisors in a
birational or finite comparison are chosen compatibly by rational
top forms. We distinguish linear equivalence $\sim$, rational
linear equivalence $\sim_{\Q}$, and numerical equivalence
$\equiv$. A pair has an effective rational boundary and a
\(\Q\)-Cartier actual adjoint. Subpairs may have negative boundary
coefficients. Log discrepancies are ordinary discrepancies plus
one; canonical and terminal singularities for zero boundary are
defined by nonnegative and positive ordinary discrepancies on
exceptional divisors. We use the usual dlt and slc conventions,
including conductor gluing in slc adjunction; see \cite{Kollar2013}.

A contraction is a surjective projective morphism $f:X\to Z$ of
normal varieties with $f_*\OO_X=\OO_Z$. A rational fibration is
understood through a relatively algebraically closed base field
inside the total function field. A log Calabi--Yau adjoint is
\(\Q\)-linearly trivial. A uniform index means a common positive
integer giving an integral linearly trivial multiple, not merely
a numerical relation. Quasi-\'{e}tale means finite surjective and
\'{e}tale outside codimension at least two. Rational connectedness
on singular spaces refers to a smooth projective model when used
in a birational argument.

Divisorial sheaves test rational sections at prime divisors. In
particular, $v\ne0$ belongs to
$H^0(X,\OO_X(\lfloor mD\rfloor))$ precisely when
$\Div(v)+mD\geq0$. Numerical positivity of a rational divisor
always concerns its rational Cartier class. Section-order and
volume limits may be taken in sufficiently divisible degrees:
powering a section preserves its normalized order. DCC and ACC
stand for the descending and ascending chain conditions.
